\section{Algebraic Equivalences}%
\label{sec:equivalences}

In addition to providing a basis to formally define the semantics of a mapping language such as RML, our mapping algebra may also be used as a foundation for implementing such languages in mapping tools.
	More precisely,
	%That is,
such tools may use our algebra
	as a form of language
to internally represent the plans that they create for converting data according to a given mapping description, and then they may use algebraic equivalences to rewrite initial mapping plans into more efficient ones. While an extensive study of such optimization opportunities is beyond the scope of this paper, in this section we show a number of such equivalences and describe how they may be relevant for a plan optimizer.

\subsection{Projection Pushing}

TODO (Olaf) some text here


The following result shows that there are cases in which a projection operator over an extend-based subexpression may be pushed completely into this subexpression. In particular, this is the case if the set of projection attributes coincides with the set of attributes mentioned by the corresponding extend expression, together with the extend attribute.
\todo{adapt}

\begin{proposition}
	\label{prop:FullProjectPushing}
	\normalfont
	Let $r = \mappingRel$ be a mapping relation,
	$\attr \in \symAttrUniverse$ be an attribute
		that is not in $\symAttrSubSet$ (i.e., $\attr \notin \symAttrSubSet$),
		%such that $\attr \notin \symAttrSubSet$,
	$\extExpr$ be an extend expression, and
	$\symProjSet \subseteq \symAttrSubSet$ be a non-empty subset of the attributes of $r$.
	If $( \fctAttrs{\extExpr} \cap \symAttrSubSet ) \subseteq \symProjSet$, then it holds that
	\begin{equation*}
		\fctProjectOpBig{\symProjSet \,\cup\, \{\attr\}\!}{ \fctExtOp{r} }
		=
		\fctExtOpBig{ \fctProjectOp{\symProjSet}{r} }
		.
	\end{equation*}
\end{proposition}

\begin{proof}
		For the proof we
		%We
	assume that $( \fctAttrs{\extExpr} \cap \symAttrSubSet ) \subseteq \symProjSet$, and we let $(\symAttrSubSet_1,\symMappingInst_1) = \fctProjectOpBig{\symProjSet \,\cup\, \{\attr\}\!}{ \fctExtOp{r} }$ and $(\symAttrSubSet_2,\symMappingInst_2) = \fctExtOpBig{ \fctProjectOp{\symProjSet}{r} }$.
	To show that $(\symAttrSubSet_1,\symMappingInst_1) = (\symAttrSubSet_2,\symMappingInst_2)$ we first note that, by Definition~\ref{def:projection_operator}, it holds that $\symAttrSubSet_1 = \symProjSet \cup \{\attr\}$, and by Definition~\ref{def:extend_operator}, $\symAttrSubSet_2 = \symProjSet \cup \{\attr\}$. Hence, $\symAttrSubSet_1 = \symAttrSubSet_2$ and, thus, it remains to show that $\symMappingInst_1 = \symMappingInst_2$.

	We first show that $\symMappingInst_1 \subseteq \symMappingInst_2$, for which we consider an arbitrary mapping tuple~$\mappingTuple \in \symMappingInst_1$ and show that $\mappingTuple \in \symMappingInst_2$.
	Given that $\mappingTuple \in \symMappingInst_1$, by Definition~\ref{def:projection_operator}, it holds that i)~$\fctDom{\mappingTuple}=\symProjSet \cup \{\attr\}$ and ii)~there exists a mapping tuple $\mappingTuple'\! \in \fctExtOp{r}$ such that $\mappingTuple(\attr') = \mappingTuple'(\attr')$ for all $\attr'\! \in \symProjSet \cup \{\attr\}$.
	Since $\mappingTuple'\! \in \fctExtOp{r}$, by Definition~\ref{def:extend_operator}, it holds that i)~$\fctDom{\mappingTuple'} = \symAttrSubSet \cup \{\attr\}$ and ii)~there exists a mapping tuple $\mappingTuple''\! \in r$ such that $\fctEval{\extExpr}{\mappingTuple''} = \mappingTuple'(\attr)$ and $\mappingTuple''(\attr') = \mappingTuple'(\attr')$ for all $\attr'\! \in \symAttrSubSet$.
	Now, let $\mappingTuple'''$ be the restriction of $\mappingTuple''$ to $\symProjSet$; i.e., $\mappingTuple'''\! = \mappingTuple''[\symProjSet]$, which, by Definition~\ref{def:projection_operator}, also means that $\fctDom{\mappingTuple'''} = \symProjSet$ and $\mappingTuple'''\! \in \fctProjectOp{\symProjSet}{r}$.
	Then, since $( \fctAttrs{\extExpr} \cap \symAttrSubSet ) \subseteq \symProjSet$ and $\mappingTuple'''(\attr') = \mappingTuple''(\attr')$ for all $\attr' \in \symProjSet$, it holds that $\fctEval{\extExpr}{\mappingTuple'''} = \fctEval{\extExpr}{\mappingTuple''}$.
	As a consequence, there exists a tuple $\mappingTuple^{(4)}\! \in \fctExtOpBig{ \fctProjectOp{\symProjSet}{r} }$ such that i)~$\fctDom{\mappingTuple^{(4)}} = \fctDom{\mappingTuple'''} \cup \{\attr\}$, ii)~$\mappingTuple^{(4)}(\attr) = \mappingTuple'(\attr)$, and iii)~$\mappingTuple^{(4)}(\attr') = \mappingTuple'''(\attr')$ for all $\attr' \in \symProjSet$.
	Since $\fctDom{\mappingTuple'''} \cup \{\attr\} = \symProjSet \cup \{\attr\} = \fctDom{\mappingTuple}$ and $\mappingTuple'(\attr) = \mappingTuple(\attr)$ and $\mappingTuple'''(\attr') = \mappingTuple''(\attr') = \mappingTuple'(\attr') = \mappingTuple(\attr')$ for all $\attr' \in \symProjSet$, we can conclude that $\mappingTuple^{(4)}\! = \mappingTuple$ and, thus, $\mappingTuple \in \fctExtOpBig{ \fctProjectOp{\symProjSet}{r} }$; i.e., $\mappingTuple \in \symMappingInst_2$.

	Now we show that $\symMappingInst_1 \supseteq \symMappingInst_2$, for which we consider a mapping tuple~$\mappingTuple \in \symMappingInst_2$ and show that $\mappingTuple \in \symMappingInst_1$.
	Given that $\mappingTuple \in \symMappingInst_2$, by Definitions~\ref{def:extend_operator} and~\ref{def:projection_operator}, it holds that i)~$\fctDom{\mappingTuple}=\{\attr\} \cup \symProjSet$ and ii)~there exists a mapping tuple $\mappingTuple'\! \in \fctProjectOp{\symProjSet}{r}$ such that $\mappingTuple(\attr) = \fctEval{\extExpr}{\mappingTuple'}$ and $\mappingTuple(\attr') = \mappingTuple'(\attr')$ for all $\attr'\! \in \symProjSet$.
	Since $\mappingTuple'\! \in \fctProjectOp{\symProjSet}{r}$, it holds that i)~$\fctDom{\mappingTuple'}=\symProjSet$ and ii)~there exists a mapping tuple $\mappingTuple''\! \in r$ such that $\fctDom{\mappingTuple''} = \symAttrSubSet$ and $\mappingTuple''(\attr') = \mappingTuple'(\attr')$ for all $\attr'\! \in \symProjSet$.
	Now, let $\mappingTuple'''$ be the mapping tuple for which it holds that i)~$\fctDom{\mappingTuple'''} = \symAttrSubSet \cup \{\attr\}$, ii)~$\mappingTuple'''(\attr) = \fctEval{\extExpr}{\mappingTuple''}$, and iii)~$\mappingTuple'''(\attr') = \mappingTuple''(\attr')$ for all $\attr'\! \in \symAttrSubSet$. Hence, by Definition~\ref{def:extend_operator}, $\mappingTuple'''\! \in \fctExtOp{r}$.
	Next, let $\mappingTuple^{(4)}$ be the restriction of $\mappingTuple'''$ to $\symProjSet \cup \{\attr\}$; i.e., $\mappingTuple^{(4)}\! = \mappingTuple'''[\symProjSet \cup \{\attr\}]$, which, by Definition~\ref{def:projection_operator}, means that $\mappingTuple^{(4)}\! \in \fctProjectOpBig{\symProjSet \,\cup\, \{\attr\}\!}{ \fctExtOp{r} }$.
	Putting everything together, we have that i)~$\fctDom{\mappingTuple^{(4)}} = \symProjSet \cup \{\attr\} = \fctDom{\mappingTuple}$, ii)~$\mappingTuple^{(4)}(\attr) = \mappingTuple'''(\attr) = \fctEval{\extExpr}{\mappingTuple''} = \fctEval{\extExpr}{\mappingTuple'} = \mappingTuple(\attr)$, and iii)~$\mappingTuple^{(4)}(\attr') = \mappingTuple'''(\attr') = \mappingTuple''(\attr') = \mappingTuple(\attr)$ for all $\attr'\! \in \symProjSet$. As a consequence, $\mappingTuple^{(4)} = \mappingTuple$ and, thus, $\mappingTuple \in \fctProjectOpBig{\symProjSet \,\cup\, \{\attr\}\!}{ \fctExtOp{r} }$; i.e., $\mappingTuple \in \symMappingInst_1$.
	\qed
\end{proof}

\newpage
\begin{proposition}
	\label{prop:ProjectSplitting}
	\normalfont
	Let $r = \mappingRel$ be a mapping relation and
	$\symProjSet \subseteq \symAttrSubSet$ be a non-empty subset of the attributes of $r$.
	For every $\symProjSet'\! \subseteq \symAttrSubSet$, it holds that
	\begin{equation*}
		\fctProjectOp{\symProjSet}{r}
		=
		\fctProjectOpBig{\symProjSet}{ \fctProjectOp{\symProjSet\cup\symProjSet'}{r} }
		.
	\end{equation*}
\end{proposition}


\begin{corollary}
	\label{prop:PartialProjectPushing}
	\normalfont
	Let $r = \mappingRel$ be a mapping relation,
	$\attr \in \symAttrUniverse$ be an attribute that is not in $\symAttrSubSet$ (i.e., $\attr \notin \symAttrSubSet$),
	$\extExpr$ be an extend expression, and
	$\symProjSet \subseteq \symAttrSubSet \cup \{\attr\}$.
	It holds that
	\begin{equation*}
		\fctProjectOpBig{\symProjSet}{ \fctExtOp{r} }
		=
		\fctProjectOpBig{\symProjSet}{ \fctExtOpBig{ \fctProjectOp{\symProjSet'\!}{r} } }
		,
	\end{equation*}
	where $\symProjSet'\! = \bigl(\symProjSet \setminus \{\attr\} \bigr) \cup \bigl(\fctAttrs{\extExpr} \cap \symAttrSubSet \bigr)$.
	%%% Note that doing "\ {a}" is important because P may contain the new
	%%% attribute a, which cannot be pushed because it is only added by the
	%%% extend operator.
	%%% Similarly, doing $\cap \symAttrSubSet$ is important because attrs(φ)
	%%% may contain attributes that are not in A, and these attributes can
	%%% also not be pushed.
\end{corollary}



\begin{proposition}
	\label{prop:ProjectionRemoval}
	\normalfont
	For every mapping relation $r = \mappingRel$, it holds that
	\begin{equation*}
		\fctProjectOp{\symAttrSubSet}{r}
		=
		r
		.
	\end{equation*}
\end{proposition}

A projection operator over a source operator may also be removed, as shown by the following result.
\begin{proposition}
	\label{prop:ProjectionIntoSource}
	\normalfont
	Let $\dataSource = \dataSourceTuple$ be a data source with
	$\sourceType=\sourceTypeTuple$,
	$\queryExpr_{root} \in \symDataAcc$,
	$\symAttrQueryMap$ be a partial function $\symAttrQueryMap \!: \symAttrUniverse  \rightarrow \symDataAcc$,
	and $\symProjSet \subseteq \fctDom{\symAttrQueryMap}$.
	It holds that
	\begin{equation*}
		\fctProjectOpBig{\symProjSet}{ \sourceOpDflt }
		=
		\sourceOp{\dataSource}{\queryExpr_{root}}{\symAttrQueryMap'}
		,
	\end{equation*}
	where $\symAttrQueryMap'$ is the restriction of $\symAttrQueryMap$ to $\symProjSet$; i.e., $\fctDom{\symAttrQueryMap'} = \symProjSet$ and $\symAttrQueryMap'(\attr) = \symAttrQueryMap(\attr)$ for all $\attr \in \fctDom{\symAttrQueryMap'}$.
\end{proposition}

\subsection{Pushing or Pulling Extend Operators}
While pushing projections would be always useful in terms of the size of the intermediate mapping relations to be processed, the decision to push an extend operator that is on top of a join-based subplan into these subplans or to pull an extend operator out of such a subplan depends on the selectivity of the join. The following result shows an equivalence
	%that, thus, may be applied in both directions by a plan optimizer.
	that a plan optimizer may use in both cases.
	%that can be used in both cases.
	
TODO: WHY?

\begin{proposition}
	\label{prop:ExtendAndJoin}
	\normalfont
	Let $r_1 = \mappingRelSpec{1}$ and $r_2 = \mappingRelSpec{2}$ be mapping relations such that $\symAttrSubSet_1 \cap \symAttrSubSet_2 = \emptyset$,
	let $\symJoinAttrPairs \subseteq \symAttrSubSet_{1} \times \symAttrSubSet_{2}$,
	let $\attr \in \symAttrUniverse$ be an attribute such that $\attr \notin \symAttrSubSet_1 \cup \symAttrSubSet_2$, and
	let $\extExpr$ be an extend expression.
	If $\fctAttrs{\extExpr} \cap \symAttrSubSet_1 = \emptyset$, then it holds that
	\begin{equation*}
		\fctExtOpBig{ \fctEquiJoinOp{r_1}{r_2} }
		=
		\fctEquiJoinOpBig{r_1}{ \fctExtOp{r_2} }
		.
	\end{equation*}
\end{proposition}	

... equivalences to move a relevant extend operator closer to a join (from either direction). The equivalence of Proposition~\ref{prop:FullProjectPushing} may be used for this purpose. Another example of such an equivalence is to swap two subsequent extend operators, as shown by the following result.

\begin{proposition}
	\label{prop:ExtendSwapping}
	\normalfont
	Let $r = \mappingRel$ be a mapping relation,
	$\attr_1,\attr_2 \in \symAttrUniverse$ be two different attributes that are not in $\symAttrSubSet$ (i.e., $\attr_1 \neq \attr_2$ and $\attr_1,\attr_2 \notin \symAttrSubSet$), and
	$\extExpr_1$ and $\extExpr_2$ be two extend expressions such that $\attr_1 \notin \fctAttrs{\extExpr_2}$ and $\attr_2 \notin \fctAttrs{\extExpr_1}$.
	It holds that
	\begin{equation*}
		\fctExtOpXBig{\attr_1}{\extExpr_1}{ \fctExtOpX{\attr_2}{\extExpr_2}{r} }
		=
		\fctExtOpXBig{\attr_2}{\extExpr_2}{ \fctExtOpX{\attr_1}{\extExpr_1}{r} }
		.
	\end{equation*}
\end{proposition}

